% Part: first-order-logic
% Chapter: natural-deduction
% Section: proving-things-quant

\documentclass[../../../include/open-logic-section]{subfiles}

\begin{document}

\olfileid[pt]{fol}{ntd}{prq}

\olsection{\usetoken{P}{derivation} com Quantificadores}

\begin{ex}
Ao lidar com quantificadores, temos de cuidar para não violar a
condição de variável própria, e às vezes isso exige que brinquemos
com a ordem de execução de certas inferências. Em geral, ajuda tentar
cuidar primeiro das regras sujeitas à condição de variável própria
(elas ficarão mais abaixo na prova finalizada).

Vamos ver como daríamos !!a{derivation} de !!{formula}
$\lexists[x][\lnot !A(x)] \lif \lnot \lforall[x][!A(x)]$.
Começando como de costume, escrevemos
\begin{prooftree}
\AxiomC{}
\UnaryInfC{$\lexists[x][\lnot !A(x)]\lif \lnot \lforall[x][!A(x)]$}
\end{prooftree}
Começamos escrevendo o que seria necessário para justificar esse
último passo usando a regra \Intro{\lif}.
\begin{prooftree}
\AxiomC{$\Discharge{\lexists[x][\lnot !A(x)]}{1}$}
\DeduceC{$\lnot \lforall[x][!A(x)]$}
\DischargeRule{\Intro{\lif}}{1}
\UnaryInfC{$\lexists[x][\lnot !A(x)]\lif \lnot \lforall[x][!A(x)]$}
\end{prooftree}
Como não há uma regra óbvia para aplicar a $\lnot \lforall[x][!A(x)]$,
prosseguiremos montando a !!{derivation} de modo que possamos usar a
regra \Elim{\lexists}. Aqui devemos prestar atenção à condição de
variável própria e escolher uma constante que não apareça em
$\lexists[x][!A(x)]$ nem em quaisquer suposições das quais ela
dependa. (Como, no entanto, nenhum !!{constant} aparece, qualquer
escolha servirá.)
\begin{prooftree}
\AxiomC{$\Discharge{\lexists[x][\lnot !A(x)]}{1}$}
\AxiomC{$\Discharge{\lnot !A(a)}{2}$}
\DeduceC{$\lnot \lforall[x][!A(x)]$}
\DischargeRule{\Elim{\lexists}}{2}
\BinaryInfC{$\lnot \lforall[x][!A(x)]$}
\DischargeRule{\Intro{\lif}}{1}
\UnaryInfC{$\lexists[x][\lnot !A(x)] \lif \lnot \lforall[x][!A(x)]$}
\end{prooftree}
Para derivar $\lnot \lforall[x][!A(x)]$, tentaremos usar a regra
\Intro{\lnot}: isso exige que derivemos uma contradição, possivelmente
usando $\lforall[x][!A(x)]$ como suposição adicional. É claro que essa
contradição pode envolver a suposição $\lnot !A(a)$, que será
descarregada pela inferência \Elim{\lexists}. Podemos montá-la da
seguinte forma:
\begin{prooftree}
\AxiomC{$\Discharge{\lexists[x][\lnot !A(x)]}{1}$}
\AxiomC{$\Discharge{\lnot !A(a)}{2}, \Discharge{\lforall[x][!A(x)]}{3}$}
\DeduceC{$\lfalse$}
\DischargeRule{\Intro{\lnot}}{3}
\UnaryInfC{$\lnot \lforall[x][!A(x)]$}
\DischargeRule{\Elim{\lexists}}{2}
\BinaryInfC{$\lnot \lforall[x][!A(x)]$}
\DischargeRule{\Intro{\lif}}{1}
\UnaryInfC{$\lexists[x][\lnot !A(x)]\lif \lnot \lforall[x][!A(x)]$}
\end{prooftree}
Parece que estamos perto de obter uma contradição. A regra mais fácil
de aplicar é a \Elim{\lforall}, que não possui condição de variável
própria. Como podemos usar qualquer termo que queiramos para
substituir o $x$ universalmente quantificado, faz mais sentido
continuar usando~$a$ para chegarmos a uma contradição.
\begin{prooftree}
\AxiomC{$\Discharge{\lexists[x][\lnot !A(x)]}{1}$}
\AxiomC{$\Discharge{\lnot !A(a)}{2}$}
\AxiomC{$\Discharge{\lforall[x][!A(x)]}{3}$}
\RightLabel{\Elim{\lforall}}
\UnaryInfC{$!A(a)$}
\RightLabel{\Elim{\lnot}}
\BinaryInfC{$\lfalse$}
\DischargeRule{\Intro{\lnot}}{3}
\UnaryInfC{$\lnot \lforall[x][!A(x)]$}
\DischargeRule{\Elim{\lexists}}{2}
\BinaryInfC{$\lnot \lforall[x][!A(x)]$}
\DischargeRule{\Intro{\lif}}{1}
\UnaryInfC{$\lexists[x][\lnot !A(x)]\lif \lnot \lforall[x][!A(x)]$}
\end{prooftree}

É importante, especialmente ao lidar com quantificadores, verificar com
cuidado neste ponto que a condição de variável própria não foi
violada. Como a única regra que aplicamos sujeita à condição de
variável própria foi \Elim{\exists}, e a variável própria~$a$ não
ocorre em nenhuma suposição da qual dependa, esta é uma
!!{derivation} correta.
\end{ex}

\begin{ex}
Às vezes podemos derivar !!a{formula} a partir de outras !!{formula}s.
Nesses casos, podemos ter suposições não descarregadas. É
importante acompanhar nossas suposições, bem como o objetivo final.

Vamos ver como daríamos !!a{derivation} de !!{formula}
$\lexists[x][!C(x,b)]$ a partir das suposições $\lexists[x][(!A(x) 
\land !B(x))]$ e $\lforall[x][(!B(x) \lif !C(x,b))]$.
Começando como de costume, escrevemos a conclusão na parte
inferior.
\begin{prooftree}
\AxiomC{}
\UnaryInfC{$\lexists[x][!C(x,b)]$}
\end{prooftree}

Temos duas premissas com as quais trabalhar. Para usar a primeira,
isto é, tentar encontrar !!a{derivation} de $\lexists[x][!C(x, b)]$
a partir de $\lexists[x][(!A(x) \land !B(x))]$, usaríamos a regra
\Elim{\lexists}. Como ela possui condição de variável própria,
aplicaremos essa regra primeiro. Obtemos o seguinte:
\begin{prooftree}
\AxiomC{$\lexists[x][(!A(x) \land !B(x))]$}
\AxiomC{$\Discharge{!A(a) \land !B(a)}{1}$}
\DeduceC{$\lexists[x][!C(x,b)]$}
\DischargeRule{\Elim{\lexists}}{1}
\BinaryInfC{$\lexists[x][!C(x,b)]$}
\end{prooftree}
As duas suposições com as quais estamos trabalhando compartilham~$!B$.
Pode ser útil neste ponto aplicar \Elim{\land} para separar~$!B(a)$.
\begin{prooftree}
\AxiomC{$\lexists[x][(!A(x) \land !B(x)])$}
\AxiomC{$\Discharge{!A(a) 
\land !B(a)}{1}$}
\RightLabel{\Elim{\land}}
\UnaryInfC{$!B(a)$}
\DeduceC{$\lexists[x][!C(x,b)]$}
\DischargeRule{\Elim{\lexists}}{1}
\BinaryInfC{$\lexists[x][!C(x,b)]$}
\end{prooftree}

A segunda suposição com a qual temos de trabalhar
é~$\lforall[x][(!B(x) \lif !C(x,b))]$. Como não há condição de
variável própria, podemos instanciar $x$ com o !!{constant}~$a$ usando
\Elim{\lforall} para obter $!B(a) \lif !C(a, b)$. Agora temos tanto
$!B(a) \lif !C(a,b)$ quanto $!B(a)$. Nosso próximo passo deve ser uma
aplicação direta da regra \Elim{\lif}.
\begin{prooftree}
\AxiomC{$\lexists[x][(!A(x) \land !B(x))]$}
\AxiomC{$\lforall[x][(!B(x) \lif !C(x,b))]$}
\RightLabel{\Elim{\lforall}}
\UnaryInfC{$!B(a) \lif !C(a,b)$}
\AxiomC{$\Discharge{!A(a) 
\land !B(a)}{1}$}
\RightLabel{\Elim{\land}}
\UnaryInfC{$!B(a)$}
\RightLabel{\Elim{\lif}}
\BinaryInfC{$!C(a,b)$}
\DeduceC{$\lexists[x][!C(x,b)]$}
\DischargeRule{\Elim{\lexists}}{1}
\insertBetweenHyps{\hspace{-5em}}
\BinaryInfC{$\lexists[x][!C(x,b)]$}
\end{prooftree}
Estamos quase lá!{} Uma aplicação de \Intro{\lexists} e atingimos
nosso objetivo.
\begin{prooftree}
\AxiomC{$\lexists[x][(!A(x) \land !B(x))]$}
\AxiomC{$\lforall[x][(!B(x) \lif !C(x,b))]$}
\RightLabel{\Elim{\lforall}}
\UnaryInfC{$!B(a) \lif !C(a,b)$}
\AxiomC{$\Discharge{!A(a) 
\land !B(a)}{1}$}
\RightLabel{\Elim{\land}}
\UnaryInfC{$!B(a)$}
\RightLabel{\Elim{\lif}}
\BinaryInfC{$!C(a,b)$}
\RightLabel{\Intro{\lexists}}
\UnaryInfC{$\lexists[x][!C(x,b)]$}
\DischargeRule{\Elim{\lexists}}{1}
\insertBetweenHyps{\hspace{-5em}}
\BinaryInfC{$\lexists[x][!C(x,b)]$}
\end{prooftree}
Como garantimos em cada etapa que as condições de variável própria
não foram violadas, podemos estar confiantes de que esta é uma
!!{derivation} correta.
\end{ex}

\begin{ex}
Dê !!a{derivation} de !!{formula}
$\lnot\lforall[x][!A(x)]$ a partir das suposições $\lforall[x][!A(x)] 
\lif \lexists[y][!B(y)]$ e $\lnot\lexists[y][!B(y)]$.
Começando como de costume, escrevemos a !!{formula} alvo na parte
inferior.
\begin{prooftree}
\AxiomC{}
\UnaryInfC{$\lnot\lforall[x][!A(x)]$}
\end{prooftree}
A última linha da !!{derivation} é uma negação, então vamos tentar
usar \Intro{\lnot}. Isso exigirá que descubramos como derivar uma
contradição.
\begin{prooftree}
\AxiomC{$\Discharge{\lforall[x][!A(x)]}{1}$}
\DeduceC{$\lfalse$}
\DischargeRule{\Intro{\lnot}}{1}
\UnaryInfC{$\lnot\lforall[x][!A(x)]$}
\end{prooftree}
Até aqui, tudo bem. Podemos usar \Elim{\lforall}, mas não é óbvio se
isso nos ajudará a alcançar nosso objetivo. Em vez disso, vamos usar
uma de nossas suposições. $\lforall[x][!A(x)] \lif \lexists[y][!B(y)]$
junto com $\lforall[x][!A(x)]$ nos permitirá usar a regra \Elim{\lif}.
\begin{prooftree}
\AxiomC{$\lforall[x][!A(x)] \lif \lexists[y][!B(y)]$}
\AxiomC{$\Discharge{\lforall[x][!A(x)]}{1}$}
\RightLabel{\Elim{\lif}}
\BinaryInfC{$\lexists[y][!B(y)]$}
\DeduceC{$\lfalse$}
\DischargeRule{\Intro{\lnot}}{1}
\UnaryInfC{$\lnot\lforall[x][!A(x)]$}
\end{prooftree}
Agora temos uma suposição final com a qual trabalhar, e parece que
isso nos ajudará a chegar a uma contradição usando \Elim{\lnot}.
\begin{prooftree}
\AxiomC{$\lnot\lexists[y][!B(y)]$}
\AxiomC{$\lforall[x][!A(x)] \lif \lexists[y][!B(y)]$}
\AxiomC{$\Discharge{\lforall[x][!A(x)]}{1}$}
\RightLabel{\Elim{\lif}}
\BinaryInfC{$\lexists[y][!B(y)]$}
\RightLabel{\Elim{\lnot}}
\BinaryInfC{$\lfalse$}
\DischargeRule{\Intro{\lnot}}{1}
\UnaryInfC{$\lnot\lforall[x][!A(x)]$}
\end{prooftree}
\end{ex}

\begin{prob}
Dê !!{derivation}s que mostrem o seguinte:
\begin{enumerate}
\item $\Proves (\lforall[x][!A(x)] \land \lforall[y][!B(y)]) \lif
\lforall[z][(!A(z) \land !B(z))]$.
\item $\Proves (\lexists[x][!A(x)] \lor \lexists[y][!B(y)]) \lif
\lexists[z][(!A(z) \lor !B(z))]$.
\item $\lforall[x][(!A(x) \lif !B)] \Proves \lexists[y][!A(y)] \lif !B$.
\item $\lforall[x][\lnot !A(x)] \Proves \lnot\lexists[x][!A(x)]$.
\item $\Proves \lnot\lexists[x][!A(x)] \lif \lforall[x][\lnot !A(x)]$.
\item $\Proves \lnot\lexists[x][\lforall[y][((!A(x,y) \lif \lnot
!A(y,y)) \land (\lnot !A(y,y) \lif !A(x,y)))]]$.
\end{enumerate}
\end{prob}

\begin{prob}
Dê !!{derivation}s que mostrem o seguinte:
\begin{enumerate}
\item $\Proves \lnot\lforall[x][!A(x)] \lif \lexists[x][\lnot!A(x)]$.
\item $(\lforall[x][!A(x)] \lif !B) \Proves \lexists[y][(!A(y) \lif !B)]$.
\item $\Proves \lexists[x][(!A(x) \lif \lforall[y][!A(y)])]$.
\end{enumerate}
(Todas exigem a regra $\FalseCl$~.)
\end{prob}

\end{document}
